\newpage
\part{Building as an Individual Developer}

Not every MCP server comes from a public registry.
Developers often write \textbf{private} or \textbf{custom} servers for local scripts, internal APIs, or personal workflows.
Each server still runs as its own process with its own transport; the \textbf{host} orchestrates them the same way as published packages.

\section{Naming Your Work}

Because \textbf{hub} usually refers to the client engine or a public marketplace, labeling a personal server collection as a ``hub'' can confuse others.

Clearer terms for a repository of custom servers include:

\begin{itemize}
  \item \textbf{MCP Monorepo}: one repository containing multiple independent MCP packages.
  \item \textbf{Custom MCP Suite}: a curated bundle of tools tailored to a specific workflow.
  \item \textbf{MCP Config or Profile}: a standardized setup shared across IDEs and clients.
\end{itemize}

\noindent
The \textbf{IDE or host is the hub} that spins up and coordinates processes; your repository is the collection of servers that hub can load.

\section{Structuring a Personal MCP Project}

A single-user setup benefits from treating custom servers as isolated packages under one repository.

\begin{itemize}
  \item Each service runs in its own \textbf{process} and communicates over its own \textbf{STDIO} or \textbf{HTTP/SSE} channel.
  \item A failure or heavy load in one server does not block the others.
  \item Shared logic (types, utilities, business rules) can live in a common library; transport and tool bindings stay in separate server packages.
  \item One Git repository and one workspace config make it easy to version, clone, and reuse the same tools across IDEs.
\end{itemize}

Use a workspace layout (npm, pnpm, or Python \lstinline[style=inline]|uv| workspaces) so dependencies do not bleed between servers.
Point each entry in your host configuration file at the correct package launch command.
